\begin{pastexam}[label={px:ust-115-ee-ds-30}]{臺聯大電機類 115 學年度 第 30 題}
（單選；2 分）

Answer \textbf{Questions 28-30} according to the following C code.
\begin{minted}[linenos, numbersep=1.5mm, xleftmargin=2em, fontsize=\small]{c}
#include <stdio.h>
#include <stdlib.h>
int counter=0;
void max_heapify(int A[],int n,int i) {
    int largest=i;
    int left=2*i+1;
    int right=2*i+2;
    counter++;
    if (left<n && A[left]>A[i])
        largest=left;
    if (right<n && A[right]>A[largest])
        largest=right;
    if (largest!=i) {
        int temp=A[i];
        A[i]=A[largest];
        A[largest]=temp;
        max_heapify(A,n,largest);
    }
}
void build_heap(int A[],int n) {
    for (int i = (n/2)-1;i>=0;i--)
        max_heapify(A,n,i);
}
int main(){
    int A[]={4,1,3,2,16,9,10,14,8,7};
    int n = sizeof(A)/sizeof(A[0]);
    build_heap(A,n);
    for (int i=0;i<n;++i)
        printf("%d,",A[i]);
    printf("\n");
    printf("counter:%d\n",counter);
    return 0;
}
\end{minted}

What is the computational complexity of \texttt{\textbf{build\_heap()}}?

(A)~$O(n)$ \qquad (B)~$O(\log n)$ \qquad (C)~$O(n^2)$ \qquad (D)~$O(n^3)$ \qquad (E)~$O(n^{1/2})$
\end{pastexam}

\begin{pastsoln}
答案是 (A)。

粗估：呼叫 \texttt{max\_heapify} 約 $n/2$ 次，每次最多下沉樹高 $\lfloor \lg n \rfloor$ 層，得到
$O(n \log n)$，但這個上界不緊。 細算（資結書定理~8.1）：高度為 $h$ 的節點最多
$\lceil n / 2^{h+1} \rceil$ 個，從它開始的 \texttt{max\_heapify} 是 $O(h)$，所以總時間
\[
\sum_{h=0}^{\lfloor \lg n \rfloor} \left\lceil \frac{n}{2^{h+1}} \right\rceil O(h)
= O\!\left(n \sum_{h=0}^{\infty} \frac{h}{2^{h}}\right) = O(2n) = O(n) .
\]
直觀地說，一半的節點是葉子（不必下沉），四分之一的節點最多下沉一層，……，只有根要下沉 $\lg n$ 層，
大部分節點都很矮。 而每個元素至少要看一次，所以是 $\Theta(n)$。

其他選項：(B)(E) 比讀一遍陣列還少，不可能；(C)(D) 都不是緊的上界（(C) 雖然數學上成立，但題目問的
是複雜度）。
\end{pastsoln}
